% =========================================================
% comandos.tex
% Colección de comandos personalizados organizados por tipo
% =========================================================
%
% Este fichero define atajos para:
% 1) Citas y notas al pie
% 2) Cabeceras
% 3) Enlaces y pies de figura
% 4) Etiquetas visuales de respuesta
% 5) Ajustes globales
% 6) Índices y glosario
%
% Recomendación:
% Mantén este archivo cargado en el preámbulo con:
% % =========================================================
% comandos.tex
% Colección de comandos personalizados organizados por tipo
% =========================================================
%
% Este fichero define atajos para:
% 1) Citas y notas al pie
% 2) Cabeceras
% 3) Enlaces y pies de figura
% 4) Etiquetas visuales de respuesta
% 5) Ajustes globales
% 6) Índices y glosario
%
% Recomendación:
% Mantén este archivo cargado en el preámbulo con:
% % =========================================================
% comandos.tex
% Colección de comandos personalizados organizados por tipo
% =========================================================
%
% Este fichero define atajos para:
% 1) Citas y notas al pie
% 2) Cabeceras
% 3) Enlaces y pies de figura
% 4) Etiquetas visuales de respuesta
% 5) Ajustes globales
% 6) Índices y glosario
%
% Recomendación:
% Mantén este archivo cargado en el preámbulo con:
% % =========================================================
% comandos.tex
% Colección de comandos personalizados organizados por tipo
% =========================================================
%
% Este fichero define atajos para:
% 1) Citas y notas al pie
% 2) Cabeceras
% 3) Enlaces y pies de figura
% 4) Etiquetas visuales de respuesta
% 5) Ajustes globales
% 6) Índices y glosario
%
% Recomendación:
% Mantén este archivo cargado en el preámbulo con:
% \input{comandos}
%
% ============================================================
% 1. Citas personalizadas en formato numérico [n]
% ============================================================
%
% Recomendado en el preámbulo:
% \usepackage[backend=biber,style=numeric,sorting=none]{biblatex}
%
% Estos comandos muestran la cita en el texto como [n], en lugar
% de superíndice, y colocan la referencia completa en nota al pie.
% Además, se añaden versiones sin nota al pie para citar solo como [n].
% ============================================================


% \bracketfootnote{contenido}
% Inserta una marca numérica [n] en el texto y añade el contenido
% correspondiente como nota al pie, también identificada con [n].
%
% Ejemplo:
% Texto citado\bracketfootnote{Referencia completa.}
\newcommand{\bracketfootnote}[1]{%
  \refstepcounter{footnote}%
  \textnormal{[\thefootnote]}%
  \begingroup
    \renewcommand{\thefootnote}{[\arabic{footnote}]}%
    \footnotetext{#1}%
  \endgroup
}


% \footcitepage{texto citado}{clave_biblatex}{pagina}
% Muestra un fragmento citado entre comillas y añade una referencia
% completa en nota al pie indicando la página concreta.
%
% Ejemplo:
% \footcitepage{La seguridad es esencial}{nist2024}{15}
%
% Resultado:
% “La seguridad es esencial” [1]
\newcommand{\footcitepage}[3]{%
  ``#1''\bracketfootnote{\fullcite{#2}, p.~#3}%
}


% \footbibliographypage{clave_biblatex}{pagina}
% Inserta una referencia completa en nota al pie indicando la página,
% sin mostrar texto citado en el cuerpo del documento.
%
% Ejemplo:
% La gestión del riesgo debe ser continua\footbibliographypage{iso27001}{12}.
\newcommand{\footbibliographypage}[2]{%
  \bracketfootnote{\fullcite{#1}, p.~#2}%
}


% \footwebfragment{clave_biblatex}{fragmento}
% Cita una fuente web en nota al pie e indica el fragmento concreto
% utilizado de esa fuente.
%
% Ejemplo:
% OWASP recomienda validar las entradas\footwebfragment{owasp2024}{Input validation}.
\newcommand{\footwebfragment}[2]{%
  \bracketfootnote{\fullcite{#1}. Fragmento: #2}%
}


% \footweb{clave_biblatex}
% Inserta una referencia web completa en nota al pie.
%
% Ejemplo:
% La guía puede consultarse en OWASP\footweb{owasp_top10_2025}.
\newcommand{\footweb}[1]{%
  \bracketfootnote{\fullcite{#1}}%
}


% \footwebcite{texto citado}{clave_biblatex}{fragmento}
% Muestra un texto citado entre comillas y añade en nota al pie
% la fuente web junto con el fragmento relevante.
%
% Ejemplo:
% \footwebcite{La autenticación multifactor reduce riesgos}{cisa2024}{MFA helps prevent unauthorized access}
\newcommand{\footwebcite}[3]{%
  ``#1''\bracketfootnote{\fullcite{#2}. Fragmento: #3}%
}


% \footcitedefinition{termino_o_definicion}{clave_biblatex}{pagina}
% Inserta una definición o término en el texto y añade la referencia
% completa en nota al pie indicando la página concreta.
%
% Ejemplo:
% \footcitedefinition{Confidencialidad}{iso27001}{8}
\newcommand{\footcitedefinition}[3]{%
  #1\bracketfootnote{\fullcite{#2}, p.~#3}%
}


% ============================================================
% Versiones sin nota al pie
% ============================================================

% \numcite{clave_biblatex}
% Inserta únicamente la cita bibliográfica numérica [n] en el texto,
% sin crear nota al pie.
%
% Ejemplo:
% La norma ISO/IEC 27001 establece controles de seguridad \numcite{iso27001}.
\newcommand{\numcite}[1]{%
  \cite{#1}%
}


% \numcitepage{clave_biblatex}{pagina}
% Inserta únicamente la cita numérica con página, sin nota al pie.
%
% Ejemplo:
% La gestión del riesgo debe ser continua \numcitepage{iso27001}{12}.
\newcommand{\numcitepage}[2]{%
  \cite[p.~#2]{#1}%
}


% \numweb{clave_biblatex}
% Alias semántico de \numcite para fuentes web.
%
% Ejemplo:
% OWASP Top 10 clasifica los riesgos principales \numweb{owasp_top10_2025}.
\newcommand{\numweb}[1]{%
  \cite{#1}%
}
% ---------------------------------------------------------
% 2. CABECERAS
% ---------------------------------------------------------
%
% Estos comandos están pensados para usar con fancyhdr.
% Permiten personalizar el contenido de la cabecera.
%
% Posiciones típicas en fancyhdr:
% - L = izquierda
% - C = centro
% - R = derecha
% - E = página par
% - O = página impar
%
% Ejemplos de posiciones:
% - LE, RO, RH, LH, etc.
%

% \nombreCapituloCabecera{posicion}{texto}
% Sirve para poner un texto en cursiva en una posición concreta
% de la cabecera.
%
% Ejemplo:
% \nombreCapituloCabecera{RH}{Capítulo 1}
\newcommand{\nombreCapituloCabecera}[2]{%
  \fancyhead[#1]{\textit{#2}}%
}

% \logocabeceraCEU{posicion}{ruta_imagen}
% Sirve para colocar un logotipo o imagen en la cabecera.
%
% Ejemplo:
% \logocabeceraCEU{LH}{Imagenes/logo-ceu.png}
%
% Nota:
% La altura de la imagen está fijada a 1 cm.
\newcommand{\logocabeceraCEU}[2]{%
  \fancyhead[#1]{\includegraphics[height=1cm]{#2}}%
}

% ---------------------------------------------------------
% 3. ENLACES, FIGURAS Y ELEMENTOS VISUALES
% ---------------------------------------------------------

% \enlace{url}{texto_visible}
% Sirve para crear un hipervínculo con color personalizado y subrayado.
%
% Ejemplo:
% \enlace{https://www.example.com}{Sitio oficial}
%
% Requisitos habituales:
% - hyperref
% - ulem o un paquete que defina \uline
% - color/bluelink definido en el documento
\newcommand{\enlace}[2]{%
  \href{#1}{\color{bluelink}\uline{#2}}%
}

% \captionazul{texto_caption}
% Sirve para crear una caption de figura o tabla con texto en color azul.
%
% Ejemplo:
% \begin{figure}
%   ...
%   \captionazul{Arquitectura general del sistema}
% \end{figure}
%
% Nota:
% También usa el mismo texto para la entrada corta del índice de figuras/tablas.
\newcommand{\captionazul}[1]{%
  \caption[#1]{\textcolor{blueColor}{#1}}%
  \vspace{0.2cm}
}


\newcommand{\subcaptionazul}[1]{%
    \subcaption{\textcolor{blueColor}{#1}}%
    \vspace{0.2cm}%
}

% ---------------------------------------------------------
% 4. ETIQUETAS VISUALES DE RESPUESTA
% ---------------------------------------------------------
%
% Estos comandos imprimen respuestas cortas en color.
% Son útiles en tablas comparativas o matrices de decisión.
%

% \permits
% Muestra “Sí” en verde.
%
% Ejemplo:
% Acceso de administrador: \permits
\newcommand{\permits}{\textcolor{green!60!black}{Sí}}

% \denies
% Muestra “No” en rojo.
%
% Ejemplo:
% Acceso de invitado: \denies
\newcommand{\denies}{\textcolor{red!70!black}{No}}

% ---------------------------------------------------------
% 5. AJUSTES GLOBALES
% ---------------------------------------------------------
%
% Estas líneas modifican comportamientos por defecto de paquetes.
%

% Elimina la línea horizontal de la cabecera definida por fancyhdr.
\renewcommand{\headrulewidth}{0pt}

% Evita que aparezcan números de página junto a las entradas del glosario.
% Suele usarse cuando no interesa mostrar en qué páginas aparece cada término.
\renewcommand*{\glossaryentrynumbers}[1]{}

% ---------------------------------------------------------
% 6. ÍNDICES Y GLOSARIO
% ---------------------------------------------------------
%
% Estos comandos generan índices y glosario con ajustes de espaciado.
%

% \crearIndiceTOC
% Genera el índice general (table of contents).
%
% Ejemplo:
% \crearIndiceTOC
\newcommand{\crearIndiceTOC}{%
    \vspace*{-2.5cm}%
    \tableofcontents%
}

% \crearIndiceTablas
% Genera el índice de tablas.
%
% Ejemplo:
% \crearIndiceTablas
\newcommand{\crearIndiceTablas}{%
    \vspace*{-2.5cm}%
    \listoftables%
}

% \crearIndiceFiguras
% Genera el índice de figuras.
%
% Ejemplo:
% \crearIndiceFiguras
\newcommand{\crearIndiceFiguras}{%
    \vspace*{-2.5cm}%
    \listoffigures%
}

% \crearglosario
% Inserta el glosario de términos, ajusta la cabecera y añade todas
% las entradas del glosario aunque no hayan sido usadas en el texto.
%
% Ejemplo:
% \crearglosario
%
% Nota:
% Este comando asume que existe el fichero:
% Glosario/glosario.tex
\newcommand{\crearglosario}{%
    \vspace*{-1cm}%
    \nombreCapituloCabecera{RH}{Glosario de términos}%

    %\glsaddall
    \vspace*{-2cm}%
    \input{Glosario/glosario}%
}
%
% ============================================================
% 1. Citas personalizadas en formato numérico [n]
% ============================================================
%
% Recomendado en el preámbulo:
% \usepackage[backend=biber,style=numeric,sorting=none]{biblatex}
%
% Estos comandos muestran la cita en el texto como [n], en lugar
% de superíndice, y colocan la referencia completa en nota al pie.
% Además, se añaden versiones sin nota al pie para citar solo como [n].
% ============================================================


% \bracketfootnote{contenido}
% Inserta una marca numérica [n] en el texto y añade el contenido
% correspondiente como nota al pie, también identificada con [n].
%
% Ejemplo:
% Texto citado\bracketfootnote{Referencia completa.}
\newcommand{\bracketfootnote}[1]{%
  \refstepcounter{footnote}%
  \textnormal{[\thefootnote]}%
  \begingroup
    \renewcommand{\thefootnote}{[\arabic{footnote}]}%
    \footnotetext{#1}%
  \endgroup
}


% \footcitepage{texto citado}{clave_biblatex}{pagina}
% Muestra un fragmento citado entre comillas y añade una referencia
% completa en nota al pie indicando la página concreta.
%
% Ejemplo:
% \footcitepage{La seguridad es esencial}{nist2024}{15}
%
% Resultado:
% “La seguridad es esencial” [1]
\newcommand{\footcitepage}[3]{%
  ``#1''\bracketfootnote{\fullcite{#2}, p.~#3}%
}


% \footbibliographypage{clave_biblatex}{pagina}
% Inserta una referencia completa en nota al pie indicando la página,
% sin mostrar texto citado en el cuerpo del documento.
%
% Ejemplo:
% La gestión del riesgo debe ser continua\footbibliographypage{iso27001}{12}.
\newcommand{\footbibliographypage}[2]{%
  \bracketfootnote{\fullcite{#1}, p.~#2}%
}


% \footwebfragment{clave_biblatex}{fragmento}
% Cita una fuente web en nota al pie e indica el fragmento concreto
% utilizado de esa fuente.
%
% Ejemplo:
% OWASP recomienda validar las entradas\footwebfragment{owasp2024}{Input validation}.
\newcommand{\footwebfragment}[2]{%
  \bracketfootnote{\fullcite{#1}. Fragmento: #2}%
}


% \footweb{clave_biblatex}
% Inserta una referencia web completa en nota al pie.
%
% Ejemplo:
% La guía puede consultarse en OWASP\footweb{owasp_top10_2025}.
\newcommand{\footweb}[1]{%
  \bracketfootnote{\fullcite{#1}}%
}


% \footwebcite{texto citado}{clave_biblatex}{fragmento}
% Muestra un texto citado entre comillas y añade en nota al pie
% la fuente web junto con el fragmento relevante.
%
% Ejemplo:
% \footwebcite{La autenticación multifactor reduce riesgos}{cisa2024}{MFA helps prevent unauthorized access}
\newcommand{\footwebcite}[3]{%
  ``#1''\bracketfootnote{\fullcite{#2}. Fragmento: #3}%
}


% \footcitedefinition{termino_o_definicion}{clave_biblatex}{pagina}
% Inserta una definición o término en el texto y añade la referencia
% completa en nota al pie indicando la página concreta.
%
% Ejemplo:
% \footcitedefinition{Confidencialidad}{iso27001}{8}
\newcommand{\footcitedefinition}[3]{%
  #1\bracketfootnote{\fullcite{#2}, p.~#3}%
}


% ============================================================
% Versiones sin nota al pie
% ============================================================

% \numcite{clave_biblatex}
% Inserta únicamente la cita bibliográfica numérica [n] en el texto,
% sin crear nota al pie.
%
% Ejemplo:
% La norma ISO/IEC 27001 establece controles de seguridad \numcite{iso27001}.
\newcommand{\numcite}[1]{%
  \cite{#1}%
}


% \numcitepage{clave_biblatex}{pagina}
% Inserta únicamente la cita numérica con página, sin nota al pie.
%
% Ejemplo:
% La gestión del riesgo debe ser continua \numcitepage{iso27001}{12}.
\newcommand{\numcitepage}[2]{%
  \cite[p.~#2]{#1}%
}


% \numweb{clave_biblatex}
% Alias semántico de \numcite para fuentes web.
%
% Ejemplo:
% OWASP Top 10 clasifica los riesgos principales \numweb{owasp_top10_2025}.
\newcommand{\numweb}[1]{%
  \cite{#1}%
}
% ---------------------------------------------------------
% 2. CABECERAS
% ---------------------------------------------------------
%
% Estos comandos están pensados para usar con fancyhdr.
% Permiten personalizar el contenido de la cabecera.
%
% Posiciones típicas en fancyhdr:
% - L = izquierda
% - C = centro
% - R = derecha
% - E = página par
% - O = página impar
%
% Ejemplos de posiciones:
% - LE, RO, RH, LH, etc.
%

% \nombreCapituloCabecera{posicion}{texto}
% Sirve para poner un texto en cursiva en una posición concreta
% de la cabecera.
%
% Ejemplo:
% \nombreCapituloCabecera{RH}{Capítulo 1}
\newcommand{\nombreCapituloCabecera}[2]{%
  \fancyhead[#1]{\textit{#2}}%
}

% \logocabeceraCEU{posicion}{ruta_imagen}
% Sirve para colocar un logotipo o imagen en la cabecera.
%
% Ejemplo:
% \logocabeceraCEU{LH}{Imagenes/logo-ceu.png}
%
% Nota:
% La altura de la imagen está fijada a 1 cm.
\newcommand{\logocabeceraCEU}[2]{%
  \fancyhead[#1]{\includegraphics[height=1cm]{#2}}%
}

% ---------------------------------------------------------
% 3. ENLACES, FIGURAS Y ELEMENTOS VISUALES
% ---------------------------------------------------------

% \enlace{url}{texto_visible}
% Sirve para crear un hipervínculo con color personalizado y subrayado.
%
% Ejemplo:
% \enlace{https://www.example.com}{Sitio oficial}
%
% Requisitos habituales:
% - hyperref
% - ulem o un paquete que defina \uline
% - color/bluelink definido en el documento
\newcommand{\enlace}[2]{%
  \href{#1}{\color{bluelink}\uline{#2}}%
}

% \captionazul{texto_caption}
% Sirve para crear una caption de figura o tabla con texto en color azul.
%
% Ejemplo:
% \begin{figure}
%   ...
%   \captionazul{Arquitectura general del sistema}
% \end{figure}
%
% Nota:
% También usa el mismo texto para la entrada corta del índice de figuras/tablas.
\newcommand{\captionazul}[1]{%
  \caption[#1]{\textcolor{blueColor}{#1}}%
  \vspace{0.2cm}
}


\newcommand{\subcaptionazul}[1]{%
    \subcaption{\textcolor{blueColor}{#1}}%
    \vspace{0.2cm}%
}

% ---------------------------------------------------------
% 4. ETIQUETAS VISUALES DE RESPUESTA
% ---------------------------------------------------------
%
% Estos comandos imprimen respuestas cortas en color.
% Son útiles en tablas comparativas o matrices de decisión.
%

% \permits
% Muestra “Sí” en verde.
%
% Ejemplo:
% Acceso de administrador: \permits
\newcommand{\permits}{\textcolor{green!60!black}{Sí}}

% \denies
% Muestra “No” en rojo.
%
% Ejemplo:
% Acceso de invitado: \denies
\newcommand{\denies}{\textcolor{red!70!black}{No}}

% ---------------------------------------------------------
% 5. AJUSTES GLOBALES
% ---------------------------------------------------------
%
% Estas líneas modifican comportamientos por defecto de paquetes.
%

% Elimina la línea horizontal de la cabecera definida por fancyhdr.
\renewcommand{\headrulewidth}{0pt}

% Evita que aparezcan números de página junto a las entradas del glosario.
% Suele usarse cuando no interesa mostrar en qué páginas aparece cada término.
\renewcommand*{\glossaryentrynumbers}[1]{}

% ---------------------------------------------------------
% 6. ÍNDICES Y GLOSARIO
% ---------------------------------------------------------
%
% Estos comandos generan índices y glosario con ajustes de espaciado.
%

% \crearIndiceTOC
% Genera el índice general (table of contents).
%
% Ejemplo:
% \crearIndiceTOC
\newcommand{\crearIndiceTOC}{%
    \vspace*{-2.5cm}%
    \tableofcontents%
}

% \crearIndiceTablas
% Genera el índice de tablas.
%
% Ejemplo:
% \crearIndiceTablas
\newcommand{\crearIndiceTablas}{%
    \vspace*{-2.5cm}%
    \listoftables%
}

% \crearIndiceFiguras
% Genera el índice de figuras.
%
% Ejemplo:
% \crearIndiceFiguras
\newcommand{\crearIndiceFiguras}{%
    \vspace*{-2.5cm}%
    \listoffigures%
}

% \crearglosario
% Inserta el glosario de términos, ajusta la cabecera y añade todas
% las entradas del glosario aunque no hayan sido usadas en el texto.
%
% Ejemplo:
% \crearglosario
%
% Nota:
% Este comando asume que existe el fichero:
% Glosario/glosario.tex
\newcommand{\crearglosario}{%
    \vspace*{-1cm}%
    \nombreCapituloCabecera{RH}{Glosario de términos}%

    %\glsaddall
    \vspace*{-2cm}%
    % =========================================================
% glosario.tex
% Definición y generación del glosario y siglas del TFM
% =========================================================
%
% Este fichero contiene:
%   - Definiciones de términos (glosario)
%   - Definiciones de siglas (acrónimos)
%   - Estructura de impresión del glosario en el documento
%
% Se utiliza junto con el paquete:
%   \usepackage[acronym]{glossaries}
%
% ---------------------------------------------------------
% 📌 ¿Cómo usar este fichero?
% ---------------------------------------------------------
%
% 1. Define nuevos términos con:
%
%    \newglossaryentry{clave}{
%        name={Nombre visible},
%        description={Definición del término}
%    }
%
% 2. Define nuevas siglas con:
%
%    \newacronym{clave}{SIGLA}{Nombre completo}
%
% 3. Usa los términos en el documento:
%
%    \gls{clave}        → muestra el término
%
% 4. Usa las siglas:
%
%    \acrshort{clave}   → forma corta (API)
%    \acrlong{clave}    → forma larga (Interfaz de Programación de Aplicaciones)
%
% 5. Genera el glosario automáticamente con:
%
%    \crearglosario
%
% ---------------------------------------------------------
% 📖 Tipos de entradas
% ---------------------------------------------------------
%
% 🔤 Términos (glosario)
%
% Ejemplo:
%   \newglossaryentry{backend}{
%       name={Backend},
%       description={Sistema de servidores y lógica}
%   }
%
% 🔠 Siglas (acrónimos)
%
% Ejemplo:
%   \newacronym{api}{API}{Interfaz de Programación de Aplicaciones}
%
% También puedes definir plural:
%
%   \newacronym[
%     plural=SIGLAS,
%     longplural={Nombre completo en plural}
%   ]{clave}{SIGLA}{Nombre completo}
%
% ---------------------------------------------------------
% ⚙️ Impresión del glosario
% ---------------------------------------------------------
%
% Este fichero incluye dos secciones:
%
% 1. Glosario de siglas
%    - Muestra únicamente acrónimos
%
% 2. Definiciones
%    - Muestra términos definidos con \newglossaryentry
%
% Ambas secciones se insertan mediante:
%
%   \input{Glosario/glosario}
%
% y son llamadas automáticamente desde:
%
%   \crearglosario
%
% ---------------------------------------------------------
% ⚠️ Reglas importantes
% ---------------------------------------------------------
%
% - Cada entrada debe tener una clave única:
%     ❌ backend duplicado
%     ✔ backend, backend_api, backend_sistema
%
% - No usar espacios en las claves:
%     ❌ "base de datos"
%     ✔ base_datos
%
% - Usa nombres descriptivos:
%     ✔ resultset
%     ❌ r1
%
% - Evita repetir definiciones en el texto:
%     usa siempre \gls{}
%
% ---------------------------------------------------------
% 💡 Notas
% ---------------------------------------------------------
%
% - El comando \glsaddall se usa para incluir todos los términos
%   aunque no aparezcan en el texto
%
% - Puedes eliminarlo si solo quieres mostrar términos usados
%
% - El formato visual está controlado por la plantilla
%
% =========================================================

% Ejemplos de entradas de glosario y siglas. Reemplaza o añade las tuyas propias.




% =========================
% Glosario de términos
% =========================

\newglossaryentry{redteam}{
    name={Red Team},
    description={Disciplina de seguridad ofensiva orientada a simular ataques reales para evaluar la capacidad de detección y respuesta de una organización}
}

\newglossaryentry{pentesting}{
    name={Pentesting},
    description={Proceso autorizado de evaluación de la seguridad de un sistema mediante técnicas de ataque controladas}
}

\newglossaryentry{pivoting}{
    name={Pivoting},
    description={Técnica que permite utilizar un sistema comprometido como punto de acceso para alcanzar otros sistemas de una red interna}
}

\newglossaryentry{lateralmovement}{
    name={Movimiento lateral},
    description={Conjunto de técnicas utilizadas por un atacante para desplazarse entre distintos sistemas de una red tras obtener acceso inicial}
}

\newglossaryentry{hostonly}{
    name={Red Host-Only},
    description={Tipo de red virtual aislada que permite la comunicación entre máquinas virtuales y el equipo anfitrión sin exponerlas a redes externas}
}

\newglossaryentry{snapshot}{
    name={Snapshot},
    description={Captura instantánea del estado completo de una máquina virtual que permite restaurarla posteriormente}
}


% =========================
% Glosario de siglas
% =========================

\newacronym{owasp}{OWASP}{Open Worldwide Application Security Project}

\newacronym{ad}{AD}{Active Directory}

\newacronym{vm}{VM}{Máquina Virtual}

\newacronym{dns}{DNS}{Domain Name System}

\newacronym{smb}{SMB}{Server Message Block}

\newacronym{winrm}{WinRM}{Windows Remote Management}

\newacronym{http}{HTTP}{HyperText Transfer Protocol}

\newacronym{https}{HTTPS}{HyperText Transfer Protocol Secure}

\newacronym{php}{PHP}{PHP Hypertext Preprocessor}

\newacronym{sql}{SQL}{Structured Query Language}

\newacronym{mysql}{MySQL}{My Structured Query Language}

\newacronym{apache}{Apache}{Apache HTTP Server}

\newacronym{ssh}{SSH}{Secure Shell}

\newacronym{nat}{NAT}{Network Address Translation}

\newacronym{lan}{LAN}{Local Area Network}

\newacronym{ip}{IP}{Internet Protocol}

\newacronym{cpu}{CPU}{Central Processing Unit}

\newacronym{ram}{RAM}{Random Access Memory}

\newacronym{xfce}{XFCE}{XForms Common Environment}

\newacronym{mitre}{MITRE ATT\&CK}{Adversarial Tactics, Techniques and Common Knowledge}

\newacronym{sqli}{SQLi}{SQL Injection}

\newacronym{xss}{XSS}{Cross-Site Scripting}

\newacronym{dvwa}{DVWA}{Damn Vulnerable Web Application}

\newacronym{cve}{CVE}{Common Vulnerabilities and Exposures}

\newacronym{cwe}{CWE}{Common Weakness Enumeration}

\newacronym{git}{Git}{Sistema de control de versiones distribuido}

\newacronym{github}{GitHub}{Plataforma de alojamiento y gestión de repositorios Git}

\newacronym{api}{API}{Interfaz de Programación de Aplicaciones}


% =========================
% Impresión del glosario
% =========================

\glsaddall

\section*{Glosario de siglas}
\label{sec:glosario-siglas}
\vspace{-1.3cm}

\printglossary[type=\acronymtype, title={}, toctitle={}]

\section*{Definiciones}
\label{sec:glosario-terminos}
\vspace{-1.3cm}

\printglossary[title={}, toctitle={}]
%
}
%
% ============================================================
% 1. Citas personalizadas en formato numérico [n]
% ============================================================
%
% Recomendado en el preámbulo:
% \usepackage[backend=biber,style=numeric,sorting=none]{biblatex}
%
% Estos comandos muestran la cita en el texto como [n], en lugar
% de superíndice, y colocan la referencia completa en nota al pie.
% Además, se añaden versiones sin nota al pie para citar solo como [n].
% ============================================================


% \bracketfootnote{contenido}
% Inserta una marca numérica [n] en el texto y añade el contenido
% correspondiente como nota al pie, también identificada con [n].
%
% Ejemplo:
% Texto citado\bracketfootnote{Referencia completa.}
\newcommand{\bracketfootnote}[1]{%
  \refstepcounter{footnote}%
  \textnormal{[\thefootnote]}%
  \begingroup
    \renewcommand{\thefootnote}{[\arabic{footnote}]}%
    \footnotetext{#1}%
  \endgroup
}


% \footcitepage{texto citado}{clave_biblatex}{pagina}
% Muestra un fragmento citado entre comillas y añade una referencia
% completa en nota al pie indicando la página concreta.
%
% Ejemplo:
% \footcitepage{La seguridad es esencial}{nist2024}{15}
%
% Resultado:
% “La seguridad es esencial” [1]
\newcommand{\footcitepage}[3]{%
  ``#1''\bracketfootnote{\fullcite{#2}, p.~#3}%
}


% \footbibliographypage{clave_biblatex}{pagina}
% Inserta una referencia completa en nota al pie indicando la página,
% sin mostrar texto citado en el cuerpo del documento.
%
% Ejemplo:
% La gestión del riesgo debe ser continua\footbibliographypage{iso27001}{12}.
\newcommand{\footbibliographypage}[2]{%
  \bracketfootnote{\fullcite{#1}, p.~#2}%
}


% \footwebfragment{clave_biblatex}{fragmento}
% Cita una fuente web en nota al pie e indica el fragmento concreto
% utilizado de esa fuente.
%
% Ejemplo:
% OWASP recomienda validar las entradas\footwebfragment{owasp2024}{Input validation}.
\newcommand{\footwebfragment}[2]{%
  \bracketfootnote{\fullcite{#1}. Fragmento: #2}%
}


% \footweb{clave_biblatex}
% Inserta una referencia web completa en nota al pie.
%
% Ejemplo:
% La guía puede consultarse en OWASP\footweb{owasp_top10_2025}.
\newcommand{\footweb}[1]{%
  \bracketfootnote{\fullcite{#1}}%
}


% \footwebcite{texto citado}{clave_biblatex}{fragmento}
% Muestra un texto citado entre comillas y añade en nota al pie
% la fuente web junto con el fragmento relevante.
%
% Ejemplo:
% \footwebcite{La autenticación multifactor reduce riesgos}{cisa2024}{MFA helps prevent unauthorized access}
\newcommand{\footwebcite}[3]{%
  ``#1''\bracketfootnote{\fullcite{#2}. Fragmento: #3}%
}


% \footcitedefinition{termino_o_definicion}{clave_biblatex}{pagina}
% Inserta una definición o término en el texto y añade la referencia
% completa en nota al pie indicando la página concreta.
%
% Ejemplo:
% \footcitedefinition{Confidencialidad}{iso27001}{8}
\newcommand{\footcitedefinition}[3]{%
  #1\bracketfootnote{\fullcite{#2}, p.~#3}%
}


% ============================================================
% Versiones sin nota al pie
% ============================================================

% \numcite{clave_biblatex}
% Inserta únicamente la cita bibliográfica numérica [n] en el texto,
% sin crear nota al pie.
%
% Ejemplo:
% La norma ISO/IEC 27001 establece controles de seguridad \numcite{iso27001}.
\newcommand{\numcite}[1]{%
  \cite{#1}%
}


% \numcitepage{clave_biblatex}{pagina}
% Inserta únicamente la cita numérica con página, sin nota al pie.
%
% Ejemplo:
% La gestión del riesgo debe ser continua \numcitepage{iso27001}{12}.
\newcommand{\numcitepage}[2]{%
  \cite[p.~#2]{#1}%
}


% \numweb{clave_biblatex}
% Alias semántico de \numcite para fuentes web.
%
% Ejemplo:
% OWASP Top 10 clasifica los riesgos principales \numweb{owasp_top10_2025}.
\newcommand{\numweb}[1]{%
  \cite{#1}%
}
% ---------------------------------------------------------
% 2. CABECERAS
% ---------------------------------------------------------
%
% Estos comandos están pensados para usar con fancyhdr.
% Permiten personalizar el contenido de la cabecera.
%
% Posiciones típicas en fancyhdr:
% - L = izquierda
% - C = centro
% - R = derecha
% - E = página par
% - O = página impar
%
% Ejemplos de posiciones:
% - LE, RO, RH, LH, etc.
%

% \nombreCapituloCabecera{posicion}{texto}
% Sirve para poner un texto en cursiva en una posición concreta
% de la cabecera.
%
% Ejemplo:
% \nombreCapituloCabecera{RH}{Capítulo 1}
\newcommand{\nombreCapituloCabecera}[2]{%
  \fancyhead[#1]{\textit{#2}}%
}

% \logocabeceraCEU{posicion}{ruta_imagen}
% Sirve para colocar un logotipo o imagen en la cabecera.
%
% Ejemplo:
% \logocabeceraCEU{LH}{Imagenes/logo-ceu.png}
%
% Nota:
% La altura de la imagen está fijada a 1 cm.
\newcommand{\logocabeceraCEU}[2]{%
  \fancyhead[#1]{\includegraphics[height=1cm]{#2}}%
}

% ---------------------------------------------------------
% 3. ENLACES, FIGURAS Y ELEMENTOS VISUALES
% ---------------------------------------------------------

% \enlace{url}{texto_visible}
% Sirve para crear un hipervínculo con color personalizado y subrayado.
%
% Ejemplo:
% \enlace{https://www.example.com}{Sitio oficial}
%
% Requisitos habituales:
% - hyperref
% - ulem o un paquete que defina \uline
% - color/bluelink definido en el documento
\newcommand{\enlace}[2]{%
  \href{#1}{\color{bluelink}\uline{#2}}%
}

% \captionazul{texto_caption}
% Sirve para crear una caption de figura o tabla con texto en color azul.
%
% Ejemplo:
% \begin{figure}
%   ...
%   \captionazul{Arquitectura general del sistema}
% \end{figure}
%
% Nota:
% También usa el mismo texto para la entrada corta del índice de figuras/tablas.
\newcommand{\captionazul}[1]{%
  \caption[#1]{\textcolor{blueColor}{#1}}%
  \vspace{0.2cm}
}


\newcommand{\subcaptionazul}[1]{%
    \subcaption{\textcolor{blueColor}{#1}}%
    \vspace{0.2cm}%
}

% ---------------------------------------------------------
% 4. ETIQUETAS VISUALES DE RESPUESTA
% ---------------------------------------------------------
%
% Estos comandos imprimen respuestas cortas en color.
% Son útiles en tablas comparativas o matrices de decisión.
%

% \permits
% Muestra “Sí” en verde.
%
% Ejemplo:
% Acceso de administrador: \permits
\newcommand{\permits}{\textcolor{green!60!black}{Sí}}

% \denies
% Muestra “No” en rojo.
%
% Ejemplo:
% Acceso de invitado: \denies
\newcommand{\denies}{\textcolor{red!70!black}{No}}

% ---------------------------------------------------------
% 5. AJUSTES GLOBALES
% ---------------------------------------------------------
%
% Estas líneas modifican comportamientos por defecto de paquetes.
%

% Elimina la línea horizontal de la cabecera definida por fancyhdr.
\renewcommand{\headrulewidth}{0pt}

% Evita que aparezcan números de página junto a las entradas del glosario.
% Suele usarse cuando no interesa mostrar en qué páginas aparece cada término.
\renewcommand*{\glossaryentrynumbers}[1]{}

% ---------------------------------------------------------
% 6. ÍNDICES Y GLOSARIO
% ---------------------------------------------------------
%
% Estos comandos generan índices y glosario con ajustes de espaciado.
%

% \crearIndiceTOC
% Genera el índice general (table of contents).
%
% Ejemplo:
% \crearIndiceTOC
\newcommand{\crearIndiceTOC}{%
    \vspace*{-2.5cm}%
    \tableofcontents%
}

% \crearIndiceTablas
% Genera el índice de tablas.
%
% Ejemplo:
% \crearIndiceTablas
\newcommand{\crearIndiceTablas}{%
    \vspace*{-2.5cm}%
    \listoftables%
}

% \crearIndiceFiguras
% Genera el índice de figuras.
%
% Ejemplo:
% \crearIndiceFiguras
\newcommand{\crearIndiceFiguras}{%
    \vspace*{-2.5cm}%
    \listoffigures%
}

% \crearglosario
% Inserta el glosario de términos, ajusta la cabecera y añade todas
% las entradas del glosario aunque no hayan sido usadas en el texto.
%
% Ejemplo:
% \crearglosario
%
% Nota:
% Este comando asume que existe el fichero:
% Glosario/glosario.tex
\newcommand{\crearglosario}{%
    \vspace*{-1cm}%
    \nombreCapituloCabecera{RH}{Glosario de términos}%

    %\glsaddall
    \vspace*{-2cm}%
    % =========================================================
% glosario.tex
% Definición y generación del glosario y siglas del TFM
% =========================================================
%
% Este fichero contiene:
%   - Definiciones de términos (glosario)
%   - Definiciones de siglas (acrónimos)
%   - Estructura de impresión del glosario en el documento
%
% Se utiliza junto con el paquete:
%   \usepackage[acronym]{glossaries}
%
% ---------------------------------------------------------
% 📌 ¿Cómo usar este fichero?
% ---------------------------------------------------------
%
% 1. Define nuevos términos con:
%
%    \newglossaryentry{clave}{
%        name={Nombre visible},
%        description={Definición del término}
%    }
%
% 2. Define nuevas siglas con:
%
%    \newacronym{clave}{SIGLA}{Nombre completo}
%
% 3. Usa los términos en el documento:
%
%    \gls{clave}        → muestra el término
%
% 4. Usa las siglas:
%
%    \acrshort{clave}   → forma corta (API)
%    \acrlong{clave}    → forma larga (Interfaz de Programación de Aplicaciones)
%
% 5. Genera el glosario automáticamente con:
%
%    \crearglosario
%
% ---------------------------------------------------------
% 📖 Tipos de entradas
% ---------------------------------------------------------
%
% 🔤 Términos (glosario)
%
% Ejemplo:
%   \newglossaryentry{backend}{
%       name={Backend},
%       description={Sistema de servidores y lógica}
%   }
%
% 🔠 Siglas (acrónimos)
%
% Ejemplo:
%   \newacronym{api}{API}{Interfaz de Programación de Aplicaciones}
%
% También puedes definir plural:
%
%   \newacronym[
%     plural=SIGLAS,
%     longplural={Nombre completo en plural}
%   ]{clave}{SIGLA}{Nombre completo}
%
% ---------------------------------------------------------
% ⚙️ Impresión del glosario
% ---------------------------------------------------------
%
% Este fichero incluye dos secciones:
%
% 1. Glosario de siglas
%    - Muestra únicamente acrónimos
%
% 2. Definiciones
%    - Muestra términos definidos con \newglossaryentry
%
% Ambas secciones se insertan mediante:
%
%   % =========================================================
% glosario.tex
% Definición y generación del glosario y siglas del TFM
% =========================================================
%
% Este fichero contiene:
%   - Definiciones de términos (glosario)
%   - Definiciones de siglas (acrónimos)
%   - Estructura de impresión del glosario en el documento
%
% Se utiliza junto con el paquete:
%   \usepackage[acronym]{glossaries}
%
% ---------------------------------------------------------
% 📌 ¿Cómo usar este fichero?
% ---------------------------------------------------------
%
% 1. Define nuevos términos con:
%
%    \newglossaryentry{clave}{
%        name={Nombre visible},
%        description={Definición del término}
%    }
%
% 2. Define nuevas siglas con:
%
%    \newacronym{clave}{SIGLA}{Nombre completo}
%
% 3. Usa los términos en el documento:
%
%    \gls{clave}        → muestra el término
%
% 4. Usa las siglas:
%
%    \acrshort{clave}   → forma corta (API)
%    \acrlong{clave}    → forma larga (Interfaz de Programación de Aplicaciones)
%
% 5. Genera el glosario automáticamente con:
%
%    \crearglosario
%
% ---------------------------------------------------------
% 📖 Tipos de entradas
% ---------------------------------------------------------
%
% 🔤 Términos (glosario)
%
% Ejemplo:
%   \newglossaryentry{backend}{
%       name={Backend},
%       description={Sistema de servidores y lógica}
%   }
%
% 🔠 Siglas (acrónimos)
%
% Ejemplo:
%   \newacronym{api}{API}{Interfaz de Programación de Aplicaciones}
%
% También puedes definir plural:
%
%   \newacronym[
%     plural=SIGLAS,
%     longplural={Nombre completo en plural}
%   ]{clave}{SIGLA}{Nombre completo}
%
% ---------------------------------------------------------
% ⚙️ Impresión del glosario
% ---------------------------------------------------------
%
% Este fichero incluye dos secciones:
%
% 1. Glosario de siglas
%    - Muestra únicamente acrónimos
%
% 2. Definiciones
%    - Muestra términos definidos con \newglossaryentry
%
% Ambas secciones se insertan mediante:
%
%   \input{Glosario/glosario}
%
% y son llamadas automáticamente desde:
%
%   \crearglosario
%
% ---------------------------------------------------------
% ⚠️ Reglas importantes
% ---------------------------------------------------------
%
% - Cada entrada debe tener una clave única:
%     ❌ backend duplicado
%     ✔ backend, backend_api, backend_sistema
%
% - No usar espacios en las claves:
%     ❌ "base de datos"
%     ✔ base_datos
%
% - Usa nombres descriptivos:
%     ✔ resultset
%     ❌ r1
%
% - Evita repetir definiciones en el texto:
%     usa siempre \gls{}
%
% ---------------------------------------------------------
% 💡 Notas
% ---------------------------------------------------------
%
% - El comando \glsaddall se usa para incluir todos los términos
%   aunque no aparezcan en el texto
%
% - Puedes eliminarlo si solo quieres mostrar términos usados
%
% - El formato visual está controlado por la plantilla
%
% =========================================================

% Ejemplos de entradas de glosario y siglas. Reemplaza o añade las tuyas propias.




% =========================
% Glosario de términos
% =========================

\newglossaryentry{redteam}{
    name={Red Team},
    description={Disciplina de seguridad ofensiva orientada a simular ataques reales para evaluar la capacidad de detección y respuesta de una organización}
}

\newglossaryentry{pentesting}{
    name={Pentesting},
    description={Proceso autorizado de evaluación de la seguridad de un sistema mediante técnicas de ataque controladas}
}

\newglossaryentry{pivoting}{
    name={Pivoting},
    description={Técnica que permite utilizar un sistema comprometido como punto de acceso para alcanzar otros sistemas de una red interna}
}

\newglossaryentry{lateralmovement}{
    name={Movimiento lateral},
    description={Conjunto de técnicas utilizadas por un atacante para desplazarse entre distintos sistemas de una red tras obtener acceso inicial}
}

\newglossaryentry{hostonly}{
    name={Red Host-Only},
    description={Tipo de red virtual aislada que permite la comunicación entre máquinas virtuales y el equipo anfitrión sin exponerlas a redes externas}
}

\newglossaryentry{snapshot}{
    name={Snapshot},
    description={Captura instantánea del estado completo de una máquina virtual que permite restaurarla posteriormente}
}


% =========================
% Glosario de siglas
% =========================

\newacronym{owasp}{OWASP}{Open Worldwide Application Security Project}

\newacronym{ad}{AD}{Active Directory}

\newacronym{vm}{VM}{Máquina Virtual}

\newacronym{dns}{DNS}{Domain Name System}

\newacronym{smb}{SMB}{Server Message Block}

\newacronym{winrm}{WinRM}{Windows Remote Management}

\newacronym{http}{HTTP}{HyperText Transfer Protocol}

\newacronym{https}{HTTPS}{HyperText Transfer Protocol Secure}

\newacronym{php}{PHP}{PHP Hypertext Preprocessor}

\newacronym{sql}{SQL}{Structured Query Language}

\newacronym{mysql}{MySQL}{My Structured Query Language}

\newacronym{apache}{Apache}{Apache HTTP Server}

\newacronym{ssh}{SSH}{Secure Shell}

\newacronym{nat}{NAT}{Network Address Translation}

\newacronym{lan}{LAN}{Local Area Network}

\newacronym{ip}{IP}{Internet Protocol}

\newacronym{cpu}{CPU}{Central Processing Unit}

\newacronym{ram}{RAM}{Random Access Memory}

\newacronym{xfce}{XFCE}{XForms Common Environment}

\newacronym{mitre}{MITRE ATT\&CK}{Adversarial Tactics, Techniques and Common Knowledge}

\newacronym{sqli}{SQLi}{SQL Injection}

\newacronym{xss}{XSS}{Cross-Site Scripting}

\newacronym{dvwa}{DVWA}{Damn Vulnerable Web Application}

\newacronym{cve}{CVE}{Common Vulnerabilities and Exposures}

\newacronym{cwe}{CWE}{Common Weakness Enumeration}

\newacronym{git}{Git}{Sistema de control de versiones distribuido}

\newacronym{github}{GitHub}{Plataforma de alojamiento y gestión de repositorios Git}

\newacronym{api}{API}{Interfaz de Programación de Aplicaciones}


% =========================
% Impresión del glosario
% =========================

\glsaddall

\section*{Glosario de siglas}
\label{sec:glosario-siglas}
\vspace{-1.3cm}

\printglossary[type=\acronymtype, title={}, toctitle={}]

\section*{Definiciones}
\label{sec:glosario-terminos}
\vspace{-1.3cm}

\printglossary[title={}, toctitle={}]

%
% y son llamadas automáticamente desde:
%
%   \crearglosario
%
% ---------------------------------------------------------
% ⚠️ Reglas importantes
% ---------------------------------------------------------
%
% - Cada entrada debe tener una clave única:
%     ❌ backend duplicado
%     ✔ backend, backend_api, backend_sistema
%
% - No usar espacios en las claves:
%     ❌ "base de datos"
%     ✔ base_datos
%
% - Usa nombres descriptivos:
%     ✔ resultset
%     ❌ r1
%
% - Evita repetir definiciones en el texto:
%     usa siempre \gls{}
%
% ---------------------------------------------------------
% 💡 Notas
% ---------------------------------------------------------
%
% - El comando \glsaddall se usa para incluir todos los términos
%   aunque no aparezcan en el texto
%
% - Puedes eliminarlo si solo quieres mostrar términos usados
%
% - El formato visual está controlado por la plantilla
%
% =========================================================

% Ejemplos de entradas de glosario y siglas. Reemplaza o añade las tuyas propias.




% =========================
% Glosario de términos
% =========================

\newglossaryentry{redteam}{
    name={Red Team},
    description={Disciplina de seguridad ofensiva orientada a simular ataques reales para evaluar la capacidad de detección y respuesta de una organización}
}

\newglossaryentry{pentesting}{
    name={Pentesting},
    description={Proceso autorizado de evaluación de la seguridad de un sistema mediante técnicas de ataque controladas}
}

\newglossaryentry{pivoting}{
    name={Pivoting},
    description={Técnica que permite utilizar un sistema comprometido como punto de acceso para alcanzar otros sistemas de una red interna}
}

\newglossaryentry{lateralmovement}{
    name={Movimiento lateral},
    description={Conjunto de técnicas utilizadas por un atacante para desplazarse entre distintos sistemas de una red tras obtener acceso inicial}
}

\newglossaryentry{hostonly}{
    name={Red Host-Only},
    description={Tipo de red virtual aislada que permite la comunicación entre máquinas virtuales y el equipo anfitrión sin exponerlas a redes externas}
}

\newglossaryentry{snapshot}{
    name={Snapshot},
    description={Captura instantánea del estado completo de una máquina virtual que permite restaurarla posteriormente}
}


% =========================
% Glosario de siglas
% =========================

\newacronym{owasp}{OWASP}{Open Worldwide Application Security Project}

\newacronym{ad}{AD}{Active Directory}

\newacronym{vm}{VM}{Máquina Virtual}

\newacronym{dns}{DNS}{Domain Name System}

\newacronym{smb}{SMB}{Server Message Block}

\newacronym{winrm}{WinRM}{Windows Remote Management}

\newacronym{http}{HTTP}{HyperText Transfer Protocol}

\newacronym{https}{HTTPS}{HyperText Transfer Protocol Secure}

\newacronym{php}{PHP}{PHP Hypertext Preprocessor}

\newacronym{sql}{SQL}{Structured Query Language}

\newacronym{mysql}{MySQL}{My Structured Query Language}

\newacronym{apache}{Apache}{Apache HTTP Server}

\newacronym{ssh}{SSH}{Secure Shell}

\newacronym{nat}{NAT}{Network Address Translation}

\newacronym{lan}{LAN}{Local Area Network}

\newacronym{ip}{IP}{Internet Protocol}

\newacronym{cpu}{CPU}{Central Processing Unit}

\newacronym{ram}{RAM}{Random Access Memory}

\newacronym{xfce}{XFCE}{XForms Common Environment}

\newacronym{mitre}{MITRE ATT\&CK}{Adversarial Tactics, Techniques and Common Knowledge}

\newacronym{sqli}{SQLi}{SQL Injection}

\newacronym{xss}{XSS}{Cross-Site Scripting}

\newacronym{dvwa}{DVWA}{Damn Vulnerable Web Application}

\newacronym{cve}{CVE}{Common Vulnerabilities and Exposures}

\newacronym{cwe}{CWE}{Common Weakness Enumeration}

\newacronym{git}{Git}{Sistema de control de versiones distribuido}

\newacronym{github}{GitHub}{Plataforma de alojamiento y gestión de repositorios Git}

\newacronym{api}{API}{Interfaz de Programación de Aplicaciones}


% =========================
% Impresión del glosario
% =========================

\glsaddall

\section*{Glosario de siglas}
\label{sec:glosario-siglas}
\vspace{-1.3cm}

\printglossary[type=\acronymtype, title={}, toctitle={}]

\section*{Definiciones}
\label{sec:glosario-terminos}
\vspace{-1.3cm}

\printglossary[title={}, toctitle={}]
%
}
%
% ============================================================
% 1. Citas personalizadas en formato numérico [n]
% ============================================================
%
% Recomendado en el preámbulo:
% \usepackage[backend=biber,style=numeric,sorting=none]{biblatex}
%
% Estos comandos muestran la cita en el texto como [n], en lugar
% de superíndice, y colocan la referencia completa en nota al pie.
% Además, se añaden versiones sin nota al pie para citar solo como [n].
% ============================================================


% \bracketfootnote{contenido}
% Inserta una marca numérica [n] en el texto y añade el contenido
% correspondiente como nota al pie, también identificada con [n].
%
% Ejemplo:
% Texto citado\bracketfootnote{Referencia completa.}
\newcommand{\bracketfootnote}[1]{%
  \refstepcounter{footnote}%
  \textnormal{[\thefootnote]}%
  \begingroup
    \renewcommand{\thefootnote}{[\arabic{footnote}]}%
    \footnotetext{#1}%
  \endgroup
}


% \footcitepage{texto citado}{clave_biblatex}{pagina}
% Muestra un fragmento citado entre comillas y añade una referencia
% completa en nota al pie indicando la página concreta.
%
% Ejemplo:
% \footcitepage{La seguridad es esencial}{nist2024}{15}
%
% Resultado:
% “La seguridad es esencial” [1]
\newcommand{\footcitepage}[3]{%
  ``#1''\bracketfootnote{\fullcite{#2}, p.~#3}%
}


% \footbibliographypage{clave_biblatex}{pagina}
% Inserta una referencia completa en nota al pie indicando la página,
% sin mostrar texto citado en el cuerpo del documento.
%
% Ejemplo:
% La gestión del riesgo debe ser continua\footbibliographypage{iso27001}{12}.
\newcommand{\footbibliographypage}[2]{%
  \bracketfootnote{\fullcite{#1}, p.~#2}%
}


% \footwebfragment{clave_biblatex}{fragmento}
% Cita una fuente web en nota al pie e indica el fragmento concreto
% utilizado de esa fuente.
%
% Ejemplo:
% OWASP recomienda validar las entradas\footwebfragment{owasp2024}{Input validation}.
\newcommand{\footwebfragment}[2]{%
  \bracketfootnote{\fullcite{#1}. Fragmento: #2}%
}


% \footweb{clave_biblatex}
% Inserta una referencia web completa en nota al pie.
%
% Ejemplo:
% La guía puede consultarse en OWASP\footweb{owasp_top10_2025}.
\newcommand{\footweb}[1]{%
  \bracketfootnote{\fullcite{#1}}%
}


% \footwebcite{texto citado}{clave_biblatex}{fragmento}
% Muestra un texto citado entre comillas y añade en nota al pie
% la fuente web junto con el fragmento relevante.
%
% Ejemplo:
% \footwebcite{La autenticación multifactor reduce riesgos}{cisa2024}{MFA helps prevent unauthorized access}
\newcommand{\footwebcite}[3]{%
  ``#1''\bracketfootnote{\fullcite{#2}. Fragmento: #3}%
}


% \footcitedefinition{termino_o_definicion}{clave_biblatex}{pagina}
% Inserta una definición o término en el texto y añade la referencia
% completa en nota al pie indicando la página concreta.
%
% Ejemplo:
% \footcitedefinition{Confidencialidad}{iso27001}{8}
\newcommand{\footcitedefinition}[3]{%
  #1\bracketfootnote{\fullcite{#2}, p.~#3}%
}


% ============================================================
% Versiones sin nota al pie
% ============================================================

% \numcite{clave_biblatex}
% Inserta únicamente la cita bibliográfica numérica [n] en el texto,
% sin crear nota al pie.
%
% Ejemplo:
% La norma ISO/IEC 27001 establece controles de seguridad \numcite{iso27001}.
\newcommand{\numcite}[1]{%
  \cite{#1}%
}


% \numcitepage{clave_biblatex}{pagina}
% Inserta únicamente la cita numérica con página, sin nota al pie.
%
% Ejemplo:
% La gestión del riesgo debe ser continua \numcitepage{iso27001}{12}.
\newcommand{\numcitepage}[2]{%
  \cite[p.~#2]{#1}%
}


% \numweb{clave_biblatex}
% Alias semántico de \numcite para fuentes web.
%
% Ejemplo:
% OWASP Top 10 clasifica los riesgos principales \numweb{owasp_top10_2025}.
\newcommand{\numweb}[1]{%
  \cite{#1}%
}
% ---------------------------------------------------------
% 2. CABECERAS
% ---------------------------------------------------------
%
% Estos comandos están pensados para usar con fancyhdr.
% Permiten personalizar el contenido de la cabecera.
%
% Posiciones típicas en fancyhdr:
% - L = izquierda
% - C = centro
% - R = derecha
% - E = página par
% - O = página impar
%
% Ejemplos de posiciones:
% - LE, RO, RH, LH, etc.
%

% \nombreCapituloCabecera{posicion}{texto}
% Sirve para poner un texto en cursiva en una posición concreta
% de la cabecera.
%
% Ejemplo:
% \nombreCapituloCabecera{RH}{Capítulo 1}
\newcommand{\nombreCapituloCabecera}[2]{%
  \fancyhead[#1]{\textit{#2}}%
}

% \logocabeceraCEU{posicion}{ruta_imagen}
% Sirve para colocar un logotipo o imagen en la cabecera.
%
% Ejemplo:
% \logocabeceraCEU{LH}{Imagenes/logo-ceu.png}
%
% Nota:
% La altura de la imagen está fijada a 1 cm.
\newcommand{\logocabeceraCEU}[2]{%
  \fancyhead[#1]{\includegraphics[height=1cm]{#2}}%
}

% ---------------------------------------------------------
% 3. ENLACES, FIGURAS Y ELEMENTOS VISUALES
% ---------------------------------------------------------

% \enlace{url}{texto_visible}
% Sirve para crear un hipervínculo con color personalizado y subrayado.
%
% Ejemplo:
% \enlace{https://www.example.com}{Sitio oficial}
%
% Requisitos habituales:
% - hyperref
% - ulem o un paquete que defina \uline
% - color/bluelink definido en el documento
\newcommand{\enlace}[2]{%
  \href{#1}{\color{bluelink}\uline{#2}}%
}

% \captionazul{texto_caption}
% Sirve para crear una caption de figura o tabla con texto en color azul.
%
% Ejemplo:
% \begin{figure}
%   ...
%   \captionazul{Arquitectura general del sistema}
% \end{figure}
%
% Nota:
% También usa el mismo texto para la entrada corta del índice de figuras/tablas.
\newcommand{\captionazul}[1]{%
  \caption[#1]{\textcolor{blueColor}{#1}}%
  \vspace{0.2cm}
}


\newcommand{\subcaptionazul}[1]{%
    \subcaption{\textcolor{blueColor}{#1}}%
    \vspace{0.2cm}%
}

% ---------------------------------------------------------
% 4. ETIQUETAS VISUALES DE RESPUESTA
% ---------------------------------------------------------
%
% Estos comandos imprimen respuestas cortas en color.
% Son útiles en tablas comparativas o matrices de decisión.
%

% \permits
% Muestra “Sí” en verde.
%
% Ejemplo:
% Acceso de administrador: \permits
\newcommand{\permits}{\textcolor{green!60!black}{Sí}}

% \denies
% Muestra “No” en rojo.
%
% Ejemplo:
% Acceso de invitado: \denies
\newcommand{\denies}{\textcolor{red!70!black}{No}}

% ---------------------------------------------------------
% 5. AJUSTES GLOBALES
% ---------------------------------------------------------
%
% Estas líneas modifican comportamientos por defecto de paquetes.
%

% Elimina la línea horizontal de la cabecera definida por fancyhdr.
\renewcommand{\headrulewidth}{0pt}

% Evita que aparezcan números de página junto a las entradas del glosario.
% Suele usarse cuando no interesa mostrar en qué páginas aparece cada término.
\renewcommand*{\glossaryentrynumbers}[1]{}

% ---------------------------------------------------------
% 6. ÍNDICES Y GLOSARIO
% ---------------------------------------------------------
%
% Estos comandos generan índices y glosario con ajustes de espaciado.
%

% \crearIndiceTOC
% Genera el índice general (table of contents).
%
% Ejemplo:
% \crearIndiceTOC
\newcommand{\crearIndiceTOC}{%
    \vspace*{-2.5cm}%
    \tableofcontents%
}

% \crearIndiceTablas
% Genera el índice de tablas.
%
% Ejemplo:
% \crearIndiceTablas
\newcommand{\crearIndiceTablas}{%
    \vspace*{-2.5cm}%
    \listoftables%
}

% \crearIndiceFiguras
% Genera el índice de figuras.
%
% Ejemplo:
% \crearIndiceFiguras
\newcommand{\crearIndiceFiguras}{%
    \vspace*{-2.5cm}%
    \listoffigures%
}

% \crearglosario
% Inserta el glosario de términos, ajusta la cabecera y añade todas
% las entradas del glosario aunque no hayan sido usadas en el texto.
%
% Ejemplo:
% \crearglosario
%
% Nota:
% Este comando asume que existe el fichero:
% Glosario/glosario.tex
\newcommand{\crearglosario}{%
    \vspace*{-1cm}%
    \nombreCapituloCabecera{RH}{Glosario de términos}%

    %\glsaddall
    \vspace*{-2cm}%
    % =========================================================
% glosario.tex
% Definición y generación del glosario y siglas del TFM
% =========================================================
%
% Este fichero contiene:
%   - Definiciones de términos (glosario)
%   - Definiciones de siglas (acrónimos)
%   - Estructura de impresión del glosario en el documento
%
% Se utiliza junto con el paquete:
%   \usepackage[acronym]{glossaries}
%
% ---------------------------------------------------------
% 📌 ¿Cómo usar este fichero?
% ---------------------------------------------------------
%
% 1. Define nuevos términos con:
%
%    \newglossaryentry{clave}{
%        name={Nombre visible},
%        description={Definición del término}
%    }
%
% 2. Define nuevas siglas con:
%
%    \newacronym{clave}{SIGLA}{Nombre completo}
%
% 3. Usa los términos en el documento:
%
%    \gls{clave}        → muestra el término
%
% 4. Usa las siglas:
%
%    \acrshort{clave}   → forma corta (API)
%    \acrlong{clave}    → forma larga (Interfaz de Programación de Aplicaciones)
%
% 5. Genera el glosario automáticamente con:
%
%    \crearglosario
%
% ---------------------------------------------------------
% 📖 Tipos de entradas
% ---------------------------------------------------------
%
% 🔤 Términos (glosario)
%
% Ejemplo:
%   \newglossaryentry{backend}{
%       name={Backend},
%       description={Sistema de servidores y lógica}
%   }
%
% 🔠 Siglas (acrónimos)
%
% Ejemplo:
%   \newacronym{api}{API}{Interfaz de Programación de Aplicaciones}
%
% También puedes definir plural:
%
%   \newacronym[
%     plural=SIGLAS,
%     longplural={Nombre completo en plural}
%   ]{clave}{SIGLA}{Nombre completo}
%
% ---------------------------------------------------------
% ⚙️ Impresión del glosario
% ---------------------------------------------------------
%
% Este fichero incluye dos secciones:
%
% 1. Glosario de siglas
%    - Muestra únicamente acrónimos
%
% 2. Definiciones
%    - Muestra términos definidos con \newglossaryentry
%
% Ambas secciones se insertan mediante:
%
%   % =========================================================
% glosario.tex
% Definición y generación del glosario y siglas del TFM
% =========================================================
%
% Este fichero contiene:
%   - Definiciones de términos (glosario)
%   - Definiciones de siglas (acrónimos)
%   - Estructura de impresión del glosario en el documento
%
% Se utiliza junto con el paquete:
%   \usepackage[acronym]{glossaries}
%
% ---------------------------------------------------------
% 📌 ¿Cómo usar este fichero?
% ---------------------------------------------------------
%
% 1. Define nuevos términos con:
%
%    \newglossaryentry{clave}{
%        name={Nombre visible},
%        description={Definición del término}
%    }
%
% 2. Define nuevas siglas con:
%
%    \newacronym{clave}{SIGLA}{Nombre completo}
%
% 3. Usa los términos en el documento:
%
%    \gls{clave}        → muestra el término
%
% 4. Usa las siglas:
%
%    \acrshort{clave}   → forma corta (API)
%    \acrlong{clave}    → forma larga (Interfaz de Programación de Aplicaciones)
%
% 5. Genera el glosario automáticamente con:
%
%    \crearglosario
%
% ---------------------------------------------------------
% 📖 Tipos de entradas
% ---------------------------------------------------------
%
% 🔤 Términos (glosario)
%
% Ejemplo:
%   \newglossaryentry{backend}{
%       name={Backend},
%       description={Sistema de servidores y lógica}
%   }
%
% 🔠 Siglas (acrónimos)
%
% Ejemplo:
%   \newacronym{api}{API}{Interfaz de Programación de Aplicaciones}
%
% También puedes definir plural:
%
%   \newacronym[
%     plural=SIGLAS,
%     longplural={Nombre completo en plural}
%   ]{clave}{SIGLA}{Nombre completo}
%
% ---------------------------------------------------------
% ⚙️ Impresión del glosario
% ---------------------------------------------------------
%
% Este fichero incluye dos secciones:
%
% 1. Glosario de siglas
%    - Muestra únicamente acrónimos
%
% 2. Definiciones
%    - Muestra términos definidos con \newglossaryentry
%
% Ambas secciones se insertan mediante:
%
%   % =========================================================
% glosario.tex
% Definición y generación del glosario y siglas del TFM
% =========================================================
%
% Este fichero contiene:
%   - Definiciones de términos (glosario)
%   - Definiciones de siglas (acrónimos)
%   - Estructura de impresión del glosario en el documento
%
% Se utiliza junto con el paquete:
%   \usepackage[acronym]{glossaries}
%
% ---------------------------------------------------------
% 📌 ¿Cómo usar este fichero?
% ---------------------------------------------------------
%
% 1. Define nuevos términos con:
%
%    \newglossaryentry{clave}{
%        name={Nombre visible},
%        description={Definición del término}
%    }
%
% 2. Define nuevas siglas con:
%
%    \newacronym{clave}{SIGLA}{Nombre completo}
%
% 3. Usa los términos en el documento:
%
%    \gls{clave}        → muestra el término
%
% 4. Usa las siglas:
%
%    \acrshort{clave}   → forma corta (API)
%    \acrlong{clave}    → forma larga (Interfaz de Programación de Aplicaciones)
%
% 5. Genera el glosario automáticamente con:
%
%    \crearglosario
%
% ---------------------------------------------------------
% 📖 Tipos de entradas
% ---------------------------------------------------------
%
% 🔤 Términos (glosario)
%
% Ejemplo:
%   \newglossaryentry{backend}{
%       name={Backend},
%       description={Sistema de servidores y lógica}
%   }
%
% 🔠 Siglas (acrónimos)
%
% Ejemplo:
%   \newacronym{api}{API}{Interfaz de Programación de Aplicaciones}
%
% También puedes definir plural:
%
%   \newacronym[
%     plural=SIGLAS,
%     longplural={Nombre completo en plural}
%   ]{clave}{SIGLA}{Nombre completo}
%
% ---------------------------------------------------------
% ⚙️ Impresión del glosario
% ---------------------------------------------------------
%
% Este fichero incluye dos secciones:
%
% 1. Glosario de siglas
%    - Muestra únicamente acrónimos
%
% 2. Definiciones
%    - Muestra términos definidos con \newglossaryentry
%
% Ambas secciones se insertan mediante:
%
%   \input{Glosario/glosario}
%
% y son llamadas automáticamente desde:
%
%   \crearglosario
%
% ---------------------------------------------------------
% ⚠️ Reglas importantes
% ---------------------------------------------------------
%
% - Cada entrada debe tener una clave única:
%     ❌ backend duplicado
%     ✔ backend, backend_api, backend_sistema
%
% - No usar espacios en las claves:
%     ❌ "base de datos"
%     ✔ base_datos
%
% - Usa nombres descriptivos:
%     ✔ resultset
%     ❌ r1
%
% - Evita repetir definiciones en el texto:
%     usa siempre \gls{}
%
% ---------------------------------------------------------
% 💡 Notas
% ---------------------------------------------------------
%
% - El comando \glsaddall se usa para incluir todos los términos
%   aunque no aparezcan en el texto
%
% - Puedes eliminarlo si solo quieres mostrar términos usados
%
% - El formato visual está controlado por la plantilla
%
% =========================================================

% Ejemplos de entradas de glosario y siglas. Reemplaza o añade las tuyas propias.




% =========================
% Glosario de términos
% =========================

\newglossaryentry{redteam}{
    name={Red Team},
    description={Disciplina de seguridad ofensiva orientada a simular ataques reales para evaluar la capacidad de detección y respuesta de una organización}
}

\newglossaryentry{pentesting}{
    name={Pentesting},
    description={Proceso autorizado de evaluación de la seguridad de un sistema mediante técnicas de ataque controladas}
}

\newglossaryentry{pivoting}{
    name={Pivoting},
    description={Técnica que permite utilizar un sistema comprometido como punto de acceso para alcanzar otros sistemas de una red interna}
}

\newglossaryentry{lateralmovement}{
    name={Movimiento lateral},
    description={Conjunto de técnicas utilizadas por un atacante para desplazarse entre distintos sistemas de una red tras obtener acceso inicial}
}

\newglossaryentry{hostonly}{
    name={Red Host-Only},
    description={Tipo de red virtual aislada que permite la comunicación entre máquinas virtuales y el equipo anfitrión sin exponerlas a redes externas}
}

\newglossaryentry{snapshot}{
    name={Snapshot},
    description={Captura instantánea del estado completo de una máquina virtual que permite restaurarla posteriormente}
}


% =========================
% Glosario de siglas
% =========================

\newacronym{owasp}{OWASP}{Open Worldwide Application Security Project}

\newacronym{ad}{AD}{Active Directory}

\newacronym{vm}{VM}{Máquina Virtual}

\newacronym{dns}{DNS}{Domain Name System}

\newacronym{smb}{SMB}{Server Message Block}

\newacronym{winrm}{WinRM}{Windows Remote Management}

\newacronym{http}{HTTP}{HyperText Transfer Protocol}

\newacronym{https}{HTTPS}{HyperText Transfer Protocol Secure}

\newacronym{php}{PHP}{PHP Hypertext Preprocessor}

\newacronym{sql}{SQL}{Structured Query Language}

\newacronym{mysql}{MySQL}{My Structured Query Language}

\newacronym{apache}{Apache}{Apache HTTP Server}

\newacronym{ssh}{SSH}{Secure Shell}

\newacronym{nat}{NAT}{Network Address Translation}

\newacronym{lan}{LAN}{Local Area Network}

\newacronym{ip}{IP}{Internet Protocol}

\newacronym{cpu}{CPU}{Central Processing Unit}

\newacronym{ram}{RAM}{Random Access Memory}

\newacronym{xfce}{XFCE}{XForms Common Environment}

\newacronym{mitre}{MITRE ATT\&CK}{Adversarial Tactics, Techniques and Common Knowledge}

\newacronym{sqli}{SQLi}{SQL Injection}

\newacronym{xss}{XSS}{Cross-Site Scripting}

\newacronym{dvwa}{DVWA}{Damn Vulnerable Web Application}

\newacronym{cve}{CVE}{Common Vulnerabilities and Exposures}

\newacronym{cwe}{CWE}{Common Weakness Enumeration}

\newacronym{git}{Git}{Sistema de control de versiones distribuido}

\newacronym{github}{GitHub}{Plataforma de alojamiento y gestión de repositorios Git}

\newacronym{api}{API}{Interfaz de Programación de Aplicaciones}


% =========================
% Impresión del glosario
% =========================

\glsaddall

\section*{Glosario de siglas}
\label{sec:glosario-siglas}
\vspace{-1.3cm}

\printglossary[type=\acronymtype, title={}, toctitle={}]

\section*{Definiciones}
\label{sec:glosario-terminos}
\vspace{-1.3cm}

\printglossary[title={}, toctitle={}]

%
% y son llamadas automáticamente desde:
%
%   \crearglosario
%
% ---------------------------------------------------------
% ⚠️ Reglas importantes
% ---------------------------------------------------------
%
% - Cada entrada debe tener una clave única:
%     ❌ backend duplicado
%     ✔ backend, backend_api, backend_sistema
%
% - No usar espacios en las claves:
%     ❌ "base de datos"
%     ✔ base_datos
%
% - Usa nombres descriptivos:
%     ✔ resultset
%     ❌ r1
%
% - Evita repetir definiciones en el texto:
%     usa siempre \gls{}
%
% ---------------------------------------------------------
% 💡 Notas
% ---------------------------------------------------------
%
% - El comando \glsaddall se usa para incluir todos los términos
%   aunque no aparezcan en el texto
%
% - Puedes eliminarlo si solo quieres mostrar términos usados
%
% - El formato visual está controlado por la plantilla
%
% =========================================================

% Ejemplos de entradas de glosario y siglas. Reemplaza o añade las tuyas propias.




% =========================
% Glosario de términos
% =========================

\newglossaryentry{redteam}{
    name={Red Team},
    description={Disciplina de seguridad ofensiva orientada a simular ataques reales para evaluar la capacidad de detección y respuesta de una organización}
}

\newglossaryentry{pentesting}{
    name={Pentesting},
    description={Proceso autorizado de evaluación de la seguridad de un sistema mediante técnicas de ataque controladas}
}

\newglossaryentry{pivoting}{
    name={Pivoting},
    description={Técnica que permite utilizar un sistema comprometido como punto de acceso para alcanzar otros sistemas de una red interna}
}

\newglossaryentry{lateralmovement}{
    name={Movimiento lateral},
    description={Conjunto de técnicas utilizadas por un atacante para desplazarse entre distintos sistemas de una red tras obtener acceso inicial}
}

\newglossaryentry{hostonly}{
    name={Red Host-Only},
    description={Tipo de red virtual aislada que permite la comunicación entre máquinas virtuales y el equipo anfitrión sin exponerlas a redes externas}
}

\newglossaryentry{snapshot}{
    name={Snapshot},
    description={Captura instantánea del estado completo de una máquina virtual que permite restaurarla posteriormente}
}


% =========================
% Glosario de siglas
% =========================

\newacronym{owasp}{OWASP}{Open Worldwide Application Security Project}

\newacronym{ad}{AD}{Active Directory}

\newacronym{vm}{VM}{Máquina Virtual}

\newacronym{dns}{DNS}{Domain Name System}

\newacronym{smb}{SMB}{Server Message Block}

\newacronym{winrm}{WinRM}{Windows Remote Management}

\newacronym{http}{HTTP}{HyperText Transfer Protocol}

\newacronym{https}{HTTPS}{HyperText Transfer Protocol Secure}

\newacronym{php}{PHP}{PHP Hypertext Preprocessor}

\newacronym{sql}{SQL}{Structured Query Language}

\newacronym{mysql}{MySQL}{My Structured Query Language}

\newacronym{apache}{Apache}{Apache HTTP Server}

\newacronym{ssh}{SSH}{Secure Shell}

\newacronym{nat}{NAT}{Network Address Translation}

\newacronym{lan}{LAN}{Local Area Network}

\newacronym{ip}{IP}{Internet Protocol}

\newacronym{cpu}{CPU}{Central Processing Unit}

\newacronym{ram}{RAM}{Random Access Memory}

\newacronym{xfce}{XFCE}{XForms Common Environment}

\newacronym{mitre}{MITRE ATT\&CK}{Adversarial Tactics, Techniques and Common Knowledge}

\newacronym{sqli}{SQLi}{SQL Injection}

\newacronym{xss}{XSS}{Cross-Site Scripting}

\newacronym{dvwa}{DVWA}{Damn Vulnerable Web Application}

\newacronym{cve}{CVE}{Common Vulnerabilities and Exposures}

\newacronym{cwe}{CWE}{Common Weakness Enumeration}

\newacronym{git}{Git}{Sistema de control de versiones distribuido}

\newacronym{github}{GitHub}{Plataforma de alojamiento y gestión de repositorios Git}

\newacronym{api}{API}{Interfaz de Programación de Aplicaciones}


% =========================
% Impresión del glosario
% =========================

\glsaddall

\section*{Glosario de siglas}
\label{sec:glosario-siglas}
\vspace{-1.3cm}

\printglossary[type=\acronymtype, title={}, toctitle={}]

\section*{Definiciones}
\label{sec:glosario-terminos}
\vspace{-1.3cm}

\printglossary[title={}, toctitle={}]

%
% y son llamadas automáticamente desde:
%
%   \crearglosario
%
% ---------------------------------------------------------
% ⚠️ Reglas importantes
% ---------------------------------------------------------
%
% - Cada entrada debe tener una clave única:
%     ❌ backend duplicado
%     ✔ backend, backend_api, backend_sistema
%
% - No usar espacios en las claves:
%     ❌ "base de datos"
%     ✔ base_datos
%
% - Usa nombres descriptivos:
%     ✔ resultset
%     ❌ r1
%
% - Evita repetir definiciones en el texto:
%     usa siempre \gls{}
%
% ---------------------------------------------------------
% 💡 Notas
% ---------------------------------------------------------
%
% - El comando \glsaddall se usa para incluir todos los términos
%   aunque no aparezcan en el texto
%
% - Puedes eliminarlo si solo quieres mostrar términos usados
%
% - El formato visual está controlado por la plantilla
%
% =========================================================

% Ejemplos de entradas de glosario y siglas. Reemplaza o añade las tuyas propias.




% =========================
% Glosario de términos
% =========================

\newglossaryentry{redteam}{
    name={Red Team},
    description={Disciplina de seguridad ofensiva orientada a simular ataques reales para evaluar la capacidad de detección y respuesta de una organización}
}

\newglossaryentry{pentesting}{
    name={Pentesting},
    description={Proceso autorizado de evaluación de la seguridad de un sistema mediante técnicas de ataque controladas}
}

\newglossaryentry{pivoting}{
    name={Pivoting},
    description={Técnica que permite utilizar un sistema comprometido como punto de acceso para alcanzar otros sistemas de una red interna}
}

\newglossaryentry{lateralmovement}{
    name={Movimiento lateral},
    description={Conjunto de técnicas utilizadas por un atacante para desplazarse entre distintos sistemas de una red tras obtener acceso inicial}
}

\newglossaryentry{hostonly}{
    name={Red Host-Only},
    description={Tipo de red virtual aislada que permite la comunicación entre máquinas virtuales y el equipo anfitrión sin exponerlas a redes externas}
}

\newglossaryentry{snapshot}{
    name={Snapshot},
    description={Captura instantánea del estado completo de una máquina virtual que permite restaurarla posteriormente}
}


% =========================
% Glosario de siglas
% =========================

\newacronym{owasp}{OWASP}{Open Worldwide Application Security Project}

\newacronym{ad}{AD}{Active Directory}

\newacronym{vm}{VM}{Máquina Virtual}

\newacronym{dns}{DNS}{Domain Name System}

\newacronym{smb}{SMB}{Server Message Block}

\newacronym{winrm}{WinRM}{Windows Remote Management}

\newacronym{http}{HTTP}{HyperText Transfer Protocol}

\newacronym{https}{HTTPS}{HyperText Transfer Protocol Secure}

\newacronym{php}{PHP}{PHP Hypertext Preprocessor}

\newacronym{sql}{SQL}{Structured Query Language}

\newacronym{mysql}{MySQL}{My Structured Query Language}

\newacronym{apache}{Apache}{Apache HTTP Server}

\newacronym{ssh}{SSH}{Secure Shell}

\newacronym{nat}{NAT}{Network Address Translation}

\newacronym{lan}{LAN}{Local Area Network}

\newacronym{ip}{IP}{Internet Protocol}

\newacronym{cpu}{CPU}{Central Processing Unit}

\newacronym{ram}{RAM}{Random Access Memory}

\newacronym{xfce}{XFCE}{XForms Common Environment}

\newacronym{mitre}{MITRE ATT\&CK}{Adversarial Tactics, Techniques and Common Knowledge}

\newacronym{sqli}{SQLi}{SQL Injection}

\newacronym{xss}{XSS}{Cross-Site Scripting}

\newacronym{dvwa}{DVWA}{Damn Vulnerable Web Application}

\newacronym{cve}{CVE}{Common Vulnerabilities and Exposures}

\newacronym{cwe}{CWE}{Common Weakness Enumeration}

\newacronym{git}{Git}{Sistema de control de versiones distribuido}

\newacronym{github}{GitHub}{Plataforma de alojamiento y gestión de repositorios Git}

\newacronym{api}{API}{Interfaz de Programación de Aplicaciones}


% =========================
% Impresión del glosario
% =========================

\glsaddall

\section*{Glosario de siglas}
\label{sec:glosario-siglas}
\vspace{-1.3cm}

\printglossary[type=\acronymtype, title={}, toctitle={}]

\section*{Definiciones}
\label{sec:glosario-terminos}
\vspace{-1.3cm}

\printglossary[title={}, toctitle={}]
%
}